\section{Experimental methodology}
\label{sec:experiment}
\subsection{Evaluation units and eligibility}
The evaluation contains six corpora organized into seven study strata: IoT-23~\cite{garcia2020iot23} (11 targets), TON\_IoT~\cite{moustafa2021ton} (9), Gotham~\cite{belarbi2025gothamdata} (7), CTU-13~\cite{garcia2014ctu} (12), N-BaIoT~\cite{nbaiot2018} (9), and UNSW IoT attack flows~\cite{hamza2019mud} G1 and G2 (5 each). These 58 units are the eligible benchmark targets, not a random sample of deployments. Each stratum uses its own compatible feature schema. Peer transfer occurs within the prepared pool for that stratum, rather than across corpora.

Reported-target membership is predefined by corpus-specific client manifests; eligibility is then checked within those manifests. A target must have at least one eligible attack type, at least 350 calibration observations, and at least four eligible peers. The configured minimum number of retained rows for an attack type is 600 by default, reduced to 400 for IoT-23, TON\_IoT, and CTU-13, and to 30 for the two UNSW strata. Eligibility therefore depends partly on the availability of evaluation attacks even though attack labels do not enter the target detector's fitting or calibration. This distinction matters when interpreting benign-only deployment claims.

The census records 769--80,000 benign training-role rows per target. Peer counts range from four in each UNSW group to eleven in CTU-13. Gotham's seven reported targets each have seven peers, indicating that reported-target and peer-pool membership must not be assumed identical. Corpus-specific preprocessing and retained sampling precede the experimental role split. Chronological ordering of retained observations is not equivalent to evaluation of an untouched continuous capture stream.

\begin{table}[tbp]
\centering
\small
\caption{Prepared target units and feature dimensions. Counts describe eligible reported targets; peer pools can be larger. Features and sampling are adapter-specific.}
\label{tab:corpora}
\begin{tabular}{llrr}
\toprule
Stratum & Evaluation unit & Targets & Features \\
\midrule
IoT-23 & Capture/scenario & 11 & 36 \\
TON\_IoT & Destination host & 9 & 29 \\
Gotham & Device capture & 7 & 37 \\
CTU-13 & Botnet scenario & 12 & 38 \\
N-BaIoT & Device directory & 9 & 115 \\
UNSW G1 & Device & 5 & 15 \\
UNSW G2 & Device & 5 & 13 \\
\bottomrule
\end{tabular}
\end{table}


Table~\ref{tab:corpora} distinguishes the evaluation units and input dimensions. TON\_IoT's selected destinations include routers, servers, and conventional hosts, so the collection should not be read as 58 consumer IoT devices. Each adapter applies a signed logarithmic feature transformation before pooled standardization. Sampling and eligibility operate on retained rows. IoT-23 caps each raw detail bucket at 200,000 benign or 60,000 malicious rows; TON\_IoT caps each host at 150,000 benign rows and each host--attack group at 60,000. IoT-23 additionally thins detail groups larger than 20,000 rows to 2\% within million-row chunks, only for files larger than 1.5 billion bytes. TON\_IoT thins host--label groups larger than 20,000 rows to 3\% within two-million-row chunks, using an intermediate 60,000-row reservoir. Both thinning counts are rounded down. CTU-13 retains at most 200,000 normal and 60,000 attack rows per scenario and discards Background traffic. Gotham retains at most 200,000 benign rows and 40,000 per attack type. N-BaIoT applies stable exact-row deduplication within files, and its source order substitutes for timestamps. The UNSW inputs use annotated minute-level flow counters and have no adapter reservoir cap.

Gotham's features require particular caution. Five features summarize all packets within each complete second, before label filtering and reservoir sampling. A retained benign row can therefore contain aggregate information about attack or unlabeled packets in that second, and later packets contribute to earlier rows' features. Role boundaries that divide a second do not separate the raw observations contributing to these aggregates. This dependency does not establish that attack labels entered model fitting, but it prevents interpreting the benchmark as immediate causal packet-stream detection.

\subsection{Role separation and fitting budgets}
The benign observations of each prepared client are ordered by timestamp or recorded source order, using a stable ordering. Split boundaries are the floors of the cumulative fractions 0.4, 0.6, and 0.8, producing training, reference, calibration, and test roles. The first $N$ training-role observations fit the local detectors. Insufficient rows cause the corresponding budget to be ineligible; observations are not borrowed from later roles to meet it.

For each eligible attack type, the ordered retained observations are divided at the midpoint. Up to the first 10,000 observations of the first half supply peer-development examples; up to the first 10,000 observations of the second half supply evaluation examples. A target's supervised model is omitted when that client becomes the monitored target. Other clients' development attacks remain available for peer training, including attack types that may also appear at the target. The experiment thus tests transfer to a target without attack labels, rather than a general guarantee of detection of attack families absent from all peer training. These chronological splits are applied separately within clients and labels; they do not impose a common wall-clock cutoff across the network.

The pooled standardizer is fitted before budget-specific local fitting, using complete benign training roles of the loaded clients. Reference and calibration roles remain separate from benign test rows in the decision pipeline. Nevertheless, this preprocessing means that the total target data available to the experiment exceeds $N$, as explained in Section~\ref{sec:problem}. Peer classifiers also use this common scaling, so target-derived benign statistics influence their input representation. No claim of fitting the entire system from 100 or 400 total target observations is made.

Four analyses avoid mixing changes in budget with changes in eligible targets. The universal primary analysis compares $N=100$ and $400$ on all 58 targets. Universal robustness adds $N=750$ on those same targets. Broad scaling compares $N=100$, $400$, and $2000$ within the 47-target cohort supporting 2000 rows. The high-resource analysis compares $N=100$, $400$, $2000$, and $20,000$ within the 15-target cohort supporting the largest budget. Budget effects are paired within their specified cohort.

\subsection{Configuration and comparison rules}
Both $\alpha=0.01$ and $0.05$ are evaluated. The reserve fraction is 0.5, the block-size parameter 250, the multiplier cap 64, and the bisection count 20. The learning seeds are 6, 7, and 8. Component seeds are derived by target identity and component domain; these runs vary learning randomness while keeping data roles fixed. The configuration's single default alpha is not the scope of the campaign: the protocol matrix explicitly includes both levels.

The primary comparator is PACE with a mean-ensemble local branch. Both methods use the same supervised peer architecture, target-relative peer evidence, peer calibration, target roles, learning-seed schedule, and evaluation cohort. The comparator thresholds the average evidence of seven local detectors at the full local allocation. The seventh detector is an autoencoder with hidden widths 20 and 10, a symmetric decoder, 30 training epochs, batch size 256, and learning rate 0.001. This comparator differs from the six-member anchor in both composition and calibration, so the primary comparison is an assessment of the complete local-branch replacement.

Two recorded baselines also retain this seven-detector local branch. The max-fusion baseline transforms peer attack probabilities relative to target reference probabilities, takes their maximum, and calibrates that fused value once. Its beyond-reference transformation differs from the margin transformation in Eq.~\eqref{eq:evidence}: it adds the logarithm of the probability-to-reference-maximum ratio to $\log(n+1)$. The Bonferroni baseline uses equally shared margin-evidence peer thresholds with multiplier fixed to one. These baselines should not be described as peer-only detectors.

Separate peer-mean and peer-maximum ablation arms fuse PACE's margin evidence while retaining the PACE local bank. These arms distinguish the effect of fused peer calibration from the probability-based recorded max-fusion baseline. The no-local-channel arm allocates the entire nominal parameter to the peer union, whereas the reported peer branch inside PACE receives only its reserved peer share. The local-only arm uses the seven-detector mean at the entire nominal parameter. Neither standalone arm is identical to a branch extracted from the final combined rule.

Additional arms remove the feature, rarity, or response family; remove the anchor; combine anchor and bank into a single weighted union; fit the response expert on margins; or reduce the local reserve to 0.25. Every arm uses the same relevant data roles and seed schedule. These are internal architectural comparisons. The campaign does not provide implementations of every contemporary federated, transfer, or calibrated intrusion detector, so literature positioning is not an empirical superiority comparison with those methods.

\subsection{Metrics and aggregation}
For target $t$, attack type $h$, and seed $s$, detection is the fraction of held-out attack observations that alert. The target detection rate averages these fractions equally across its attack types and then across the three seeds. The pooled detection rate averages targets equally. Large captures therefore do not dominate through their row counts, although strata with more targets have greater pooled weight. False-alert rate is the fraction of benign test observations that alert, averaged over seeds within target and then over targets for its pooled mean.

The primary paired difference subtracts the comparator's target detection from PACE's. Intervals use 10,000 bootstrap replicates with seed 314159 and confidence level 0.95. Targets are sampled with replacement within corpus groups identified by the dataset enum, preserving each group's size, and their contributions are weighted by target count. The interval endpoints are linear empirical quantiles of the bootstrap means. The two UNSW study strata belong to one corpus group for this resampling. Seeds and attack observations are not resampled as independent targets.

This procedure conditions on the observed corpora and the seed-averaged target effects. It does not account fully for shared peer models, common captures, cohort selection, or dependence between target units. Intervals are descriptive evidence about this benchmark; they are not uncertainty bounds for arbitrary new IoT populations. Multiple ablations and subgroup analyses are exploratory, without a claim of simultaneous family-wise coverage.

Sweep-adjusted primary comparisons use multipliers $\{0.55,0.75,1,1.4\}$ at $\alpha=0.01$ and $\{0.4,0.75,1,1.4\}$ at $\alpha=0.05$, scaling both branch allocation parameters. PACE's complete rule, including its resolution condition, is evaluated at each stored point. Primary reporting interpolates within the realized-rate sweep and clamps outside its endpoints, separately for each seed. This procedure approximates a comparator operating point; it does not guarantee exact matching. The false-alert distribution reports quantiles over target seed means, while acceptance additionally examines exceedance among target--seed cells. These denominators are distinct.

\subsection{Internal protocol checks and reproducibility}
The primary engineering acceptance rule permits no target loss exceeding 0.05 and permits at most 3\% of targets to lose more than 0.01. It also requires positive lower bounds for nominal and sweep-adjusted pooled gains; limits mean false-alert increases to 0.004 or 0.002 at the two alpha levels; limits the increase in target--seed exceedance above $2\alpha$ to 0.03; and checks excluded-TON\_IoT, learning-seed, and leave-one-stratum-out summaries. These conditions describe the repository's decision rule. They do not constitute a formal non-inferiority hypothesis test, a registered clinical-style margin analysis, or external validation.

The outcome-independent power table is a planning calculation over assumed paired standard deviations and a meaningful effect of 0.02, using level 0.05 and power 0.8. It is not an observed-power validation of the results. At 58 targets, for example, assumed standard deviations of 0.05 and 0.10 produce minimum detectable effects of 0.018393 and 0.036787. The conclusions rely on the actual paired outcomes and their limitations rather than on selecting a favorable planning assumption.

All numerical summaries reported here are taken from the existing campaign tables. No models were retrained and no scientific campaign was rerun during manuscript preparation. Held-out learning seeds assess the fixed design within the reported cohort. Provenance checks establish artifact identity and consistency; they do not extend the evaluation to other target populations.




